%%%%%%%%%%%%%%%%%%%%%%%%%%%%%%%%%%%%%%%%%%%%%%%%%%%%%%%%%%%%%%%%%%%%%%
% Lớp 10 | Phương pháp tọa độ trong mặt phẳng (𝑂𝑥𝑦) | Đường tròn trong mặt phẳng tọa độ
% Dạng 0H9?4-6: Bài toán dùng cho tam giác, tứ giác
% Số câu: 2  (sinh tự động bởi tach_cau.py; từ nay thầy sửa trực tiếp file này)
%%%%%%%%%%%%%%%%%%%%%%%%%%%%%%%%%%%%%%%%%%%%%%%%%%%%%%%%%%%%%%%%%%%%%%

%%%%%%%%%%%%%%%% ĐÚNG/SAI %%%%%%%%%%%%%%%%

% nguồn: MaHoaNH_DungSai_Lop10
\begin{ex}%[0H9N4-6]%[ma:aa2e2db8]
	Cho tam giác $ABC$ vuông  tại $A$, có $A(0;2)$, $B(-1;-1)$, $C(6;0)$. Các khẳng định sau đây đúng hay sai?
	\choiceTF 
		{Đường tròn tâm $B$ tiếp xúc với đường thẳng $AC$ tại  điểm $C$}
		{\True Đường tròn tâm  $C$ tiếp xúc với đường thẳng $AB$ có bán kính $R=2\sqrt{10}$}
		{\True Đường tròn ngoại tiếp tam giác $ABC$  có tâm là điểm $I\left(\dfrac{5}{2};-\dfrac{1}{2}\right)$}
		{\True Đường tròn ngoại tiếp tam giác $ABC$ có đường kính bằng $5\sqrt{2}$} 
	 \loigiai{
		 Tam giác $ABC$ vuông tại $A$ có $AB=\sqrt{(-1)^2+(-1-2)^2}=\sqrt{10}$, $AC=\sqrt{6^2+2^2}=2\sqrt{10}$ và $BC=\sqrt{(6+1)^2+1^2}=5\sqrt{2}$.
		\begin{itemchoice}
			\itemch  Do $BA\perp AC$ tại $A$ nên đường tròn tâm $B$ tiếp xúc với đường thẳng $AC$ tại điểm $A$.
			\itemch Đường tròn tâm $C$ tiếp xúc với đường thẳng  $AB$ tại $A$ nên có bán kính bằng $CA=2\sqrt{10}$.
			\itemch Đường tròn ngoại tiếp tam giác  $ABC$ có tâm $I$ là trung điểm của $BC$, do đó $I\left(\dfrac{5}{2};-\dfrac{1}{2}\right)$.
			\itemch Đường tròn ngoại tiếp tam  giác $ABC$ đường kính $BC=5\sqrt{2}$.
		\end{itemchoice}
	 }
\end{ex}

% nguồn: MaHoaNH_DungSai_Lop10
\begin{ex}%[0H9H4-6]%[ma:6ac69f60]
	Trong mặt phẳng tọa độ $Oxy$, cho điểm $A(2;-4)$, đường thẳng $\Delta\colon \heva{&x=-3+2t\\ &y=1+t}$ và đường tròn $(C)\colon x^2+y^2-2x-8y-8=0$. Gọi $P$, $Q$ là các giao điểm của  $\Delta$ và $(C)$ trong đó điểm $P$ có hoành độ âm.
	\choiceTF 
		{\True Phương trình tổng quát của  đường thẳng $\Delta$ là $x-2y+5=0$}
		{Đường tròn $(T)$ tâm $A$ và tiếp xúc với đường thẳng $\Delta$ có phương trình là $(x-2)^2+(y+4)^2=9$} 
		{Đường trung trực của $PQ$ có phương trình là  $2x+y-2=0$}
		{\True Điểm $M$ thuộc $\Delta$ sao cho tam giác $AMP$ vuông  tại $M$ có hoành độ bằng $-1$}
	\loigiai{ 
		Đường thẳng $\Delta$ đi qua điểm $M_0(-3;1)$ và có véc-tơ chỉ  phương $\overrightarrow{u}=(2;1)$.\\
		 Đường tròn $(C)\colon (x-1)^2+(y-4)^2=25$ có tâm $I(1;4)$, bán kính $R=5$.
		\begin{itemchoice}
			\itemch Đường  thẳng $\Delta$ có véc-tơ pháp tuyển $\overrightarrow{n}=(1;-2)$.\\
			Phương trình tổng quát của  $\Delta$ là
			$$(x+3)-2(y-1)=0 \Leftrightarrow  x-2y+5=0.$$
			\itemch Đường tròn $(T)$ tâm $A$ và tiếp  xúc với đường thẳng $\Delta$ có bán kính là
			$$R_1=\dfrac{|2+8+5|}{\sqrt{1^2+(-2)^2}}  = 3\sqrt{5}.$$
			Phương trình đường tròn $(T)\colon  (x-2)^2+(y+4)^2=45$.
			\itemch Ta có $PQ$ là dây cung của $(C)$ nên đường trung trực của $PQ$ là đường thẳng $d$ đi qua tâm $I(1;4)$ của $(C)$ và vuông góc với $\Delta$, $d$ nhận véc-tơ $\overrightarrow{u}=(2;1)$ làm véc-tơ pháp  tuyến.\\
			Phương trình tổng  quát của $d\colon 2(x-1)+(y-4)=0 \Leftrightarrow 2x+y-6=0$.
			\itemch Điểm $P$ thuộc $\Delta$ và có hoành độ âm nên $P(-3+2t;1+t)$  ($t<\dfrac{3}{2}$).\\
			Do $P\in (C)$ nên $(2t-4)^2+(t-3)^2=25 \Leftrightarrow 5t^2-22t=0 \Leftrightarrow \hoac{&t=0 \text{ (thỏa mãn)}\\&t=\dfrac{22}{5} \text{ (loại)}}$, suy ra  $P(-3;1)$.\\
			Điểm $M\in \Delta$ nên $M(-3+2m;1+m)$ ($M$ khác $P$), $\overrightarrow{AM}=(2m-5;m+5)$, $\overrightarrow{PM}=(2m;m)$.\\ 
			Tam giác $AMP$ vuông tại $M$ nên $AM\perp PM \Leftrightarrow  \overrightarrow{AM}\cdot \overrightarrow{PM}=0$, do đó
			$$(2m-5)\cdot  2m+(m+5)\cdot m =0 \Leftrightarrow m(4m-10+m+5)=0 \Leftrightarrow \hoac{&m=0\\&m=1.}$$
			Do điểm $M$  khác điểm $P$ nên $M(-1;2)$.
		\end{itemchoice}
	 }
\end{ex}
